% Web-source links in the requested markdown form:
% [ar5iv.labs.arxiv.org](https://ar5iv.labs.arxiv.org/html/1602.07378)
% [aimsciences.org](https://www.aimsciences.org/article/doi/10.3934/jmd.2017016)
% [ar5iv.labs.arxiv.org](https://ar5iv.labs.arxiv.org/html/1606.00723)
% [math.ksu.edu](https://www.math.ksu.edu/~tanyaf/FKS.pdf)

\documentclass[11pt]{article}

\usepackage{amsmath,amssymb,amsthm}
\usepackage[margin=1in]{geometry}
\usepackage{hyperref}

\newtheorem{theorem}{Theorem}[section]
\newtheorem{lemma}[theorem]{Lemma}
\newtheorem{proposition}[theorem]{Proposition}
\newtheorem{corollary}[theorem]{Corollary}
\theoremstyle{definition}
\newtheorem{definition}[theorem]{Definition}
\theoremstyle{remark}
\newtheorem{remark}[theorem]{Remark}

\newcommand{\TT}{\mathcal T}
\newcommand{\Mod}{\operatorname{Mod}}
\newcommand{\PMod}{\operatorname{PMod}}
\newcommand{\Def}{\operatorname{Def}}
\newcommand{\Stab}{\operatorname{Stab}}
\newcommand{\Per}{\operatorname{Per}}

\title{A Dehn Twist with an Infinite Orbit on a Deformation Space}
\date{}

\begin{document}

\maketitle

\section{Statement of the result}

Let \(f:\mathbb P^1\to \mathbb P^1\) be a quadratic rational map with critical
points \(c_1,c_2\).  Assume that
\[
c_1\longmapsto f(c_1)\longmapsto f^{\circ2}(c_1)
 \longmapsto f^{\circ3}(c_1)\longmapsto c_1
\]
is a cycle of exact period \(4\), and that the other critical value
\[
v:=f(c_2)
\]
is not in this cycle.  Thus \(f\) represents a point of
\(\Per_4(0)^*\) in the usual notation.

Put
\[
a_j:=f^{\circ j}(c_1),\qquad j\in\mathbb Z/4\mathbb Z,
\]
and
\[
A:=\{a_0,a_1,a_2,a_3\},\qquad B:=A\cup\{v\}.
\]
Since \(f(A)=A\) and \(V(f)=\{f(c_1),v\}\subset B\), the deformation space
\[
\Def_B^A(f)
 :=\{\tau\in\TT(\mathbb P^1,B):f^*(\tau)=i_*(\tau)\}
\]
is defined, where \(i_*:\TT(B)\to\TT(A)\) forgets \(v\).  The local theory of
deformation spaces identifies it as a smooth one-dimensional analytic
submanifold of the two-dimensional Teichm\"uller space \(\TT(B)\); see
\cite[Theorem~1.2]{KL-arxiv}.  In this example it is nonempty: the tautological
marking of the given rational map belongs to it.

The purpose of this note is to describe an explicit curve whose squared Dehn
twist stabilizes \(\Def_B^A(f)\), and whose powers give an infinite orbit in
\(\pi_0(\Def_B^A(f))\).

\begin{definition}[The curve]
Choose an embedded closed disk \(\Delta\subset\mathbb P^1\) such that
\[
\Delta\cap B=\{c_1,v\},
\]
and let
\[
\gamma:=\partial\Delta.
\]
Thus \(\gamma\) separates
\[
\{c_1,v\}
\quad\text{from}\quad
\{f(c_1),f^{\circ2}(c_1),f^{\circ3}(c_1)\}.
\]
Both sides contain at least two points of \(B\), so \(\gamma\) is essential in
\(\mathbb P^1\setminus B\).
\end{definition}

Equivalently, after a M\"obius change of coordinates we may normalize
\[
c_1=0,\qquad f(c_1)=1,\qquad f^{\circ2}(c_1)=\infty,
\qquad f^{\circ3}(c_1)=\rho.
\]
Then
\[
B=\{0,1,\infty,\rho,v\},
\]
and \(\gamma\) is any simple closed curve separating
\[
\{0,v\}\quad\text{from}\quad\{1,\infty,\rho\}.
\]

Let \(T_\gamma\in\PMod(\mathbb P^1,B)\) denote the positive Dehn twist about
\(\gamma\).  The main result is the following.

\begin{theorem}\label{T:main}
Let \(f,A,B,\gamma\) be as above.
\begin{enumerate}
\item The squared twist \(T_\gamma^2\) preserves the deformation space:
\[
T_\gamma^2\cdot\Def_B^A(f)=\Def_B^A(f).
\]
More precisely, for every \(\tau\in\TT(B)\),
\[
f^*(T_\gamma^2\cdot\tau)=f^*(\tau),
\qquad
i_*(T_\gamma^2\cdot\tau)=i_*(\tau).
\]

\item If \(C_0\) is a connected component of \(\Def_B^A(f)\), then the
components
\[
T_\gamma^{2n}\cdot C_0,\qquad n\in\mathbb Z,
\]
are pairwise distinct.  Consequently, the action of \(T_\gamma^2\) on
\(\pi_0(\Def_B^A(f))\) has an infinite, indeed free, orbit.
\end{enumerate}
\end{theorem}

The first assertion is a direct annular lifting computation.  The second is the
stronger component-stabilizer statement proved in the course of the proof of
the disconnectedness theorem of Koch--Lodge; see \cite[proof of
Theorem~1.3]{KL-arxiv} and \cite{KL-published}.

\section{The annular lift}

Write
\[
\Delta\cap B=\{a_0,v\},\qquad a_0=c_1.
\]
The only critical value of \(f\) contained in \(\Delta\) is \(v\): the other
critical value is \(f(c_1)=a_1\), which lies outside \(\Delta\).

\begin{lemma}\label{L:preimage-disk}
The preimage \(f^{-1}(\Delta)\) is a single closed disk \(\widetilde\Delta\),
and
\[
f|_{\widetilde\Delta}:\widetilde\Delta\longrightarrow\Delta
\]
is a degree-two disk map branched over \(v\).  In particular,
\(\beta:=\partial\widetilde\Delta\) is connected and
\(f|_\beta:\beta\to\gamma\) has degree \(2\).

Moreover,
\[
\widetilde\Delta\cap A=\{a_3\}.
\]
\end{lemma}

\begin{proof}
Since \(v\) is a simple critical value of a degree-two branched cover, its full
preimage is the critical point \(c_2\), counted with local degree \(2\).  Every
component of \(f^{-1}(\Delta)\) therefore contains \(c_2\), so there is only
one such component.  Riemann--Hurwitz shows that it is a disk and that its
boundary \(\beta\) maps to \(\gamma\) with degree \(2\).

If \(x\in A\cap\widetilde\Delta\), then \(f(x)\in\Delta\cap B=\{a_0,v\}\).
But \(f(A)=A\) and \(v\notin A\), so \(f(x)=a_0\).  On the four-cycle
\(a_0\mapsto a_1\mapsto a_2\mapsto a_3\mapsto a_0\), the unique preimage of
\(a_0\) lying in \(A\) is \(a_3\).  Hence
\(\widetilde\Delta\cap A=\{a_3\}\).
\end{proof}

\begin{lemma}\label{L:twist-lift}
The square \(T_\gamma^2\) lifts under \(f\) to a Dehn twist about \(\beta\),
up to sign.  After the points in \(f^{-1}(B)\setminus A\) are erased, this
lift induces the identity on \(\TT(A)\).
\end{lemma}

\begin{proof}
Choose an annular neighborhood \(N\) of \(\gamma\).  Its preimage is an
annular neighborhood \(\widetilde N\) of \(\beta\), and
\(\widetilde N\to N\) has degree \(2\) on the core curves.  The standard Dehn
twist lifting formula therefore says that \(T_\gamma^2\) lifts to
\(T_\beta^{\pm1}\).  More generally, if the preimage components of a curve map
with degrees \(d_j\), and \(m\) is divisible by every \(d_j\), then
\[
f^*(T_\gamma^m)=\prod_j T_{\beta_j}^{\pm m/d_j};
\]
see the pullback computation in \cite{FKS-arxiv,FKS-pdf}.

By Lemma~\ref{L:preimage-disk}, the disk bounded by \(\beta\) contains only
the point \(a_3\) of \(A\).  Hence \(\beta\) is peripheral in
\(\mathbb P^1\setminus A\).  Its Dehn twist is trivial in
\(\PMod(\mathbb P^1,A)\), and therefore acts trivially on \(\TT(A)\).
\end{proof}

\begin{lemma}\label{L:forget-trivial}
Forgetting \(v\) sends \(\gamma\) to a peripheral curve in
\((\mathbb P^1,A)\).  Consequently,
\[
i_*(T_\gamma^{2}\cdot\tau)=i_*(\tau)
\]
for every \(\tau\in\TT(B)\).
\end{lemma}

\begin{proof}
After deleting \(v\), the disk \(\Delta\) contains only \(c_1\) among the
points of \(A\).  Thus \(\gamma\) bounds a once-punctured disk in
\(\mathbb P^1\setminus A\), and its Dehn twist is trivial in
\(\PMod(\mathbb P^1,A)\).
\end{proof}

Equivalently, in the Birman point-pushing interpretation, \(T_\gamma^2\)
represents the point-push of \(v\) around \(c_1\), up to twists about
peripheral curves.  Those peripheral twists are trivial in the pure mapping
class group.

\section{Stabilization of the deformation space}

We use the right action of the pure mapping class group on Teichm\"uller
space.

\begin{proposition}\label{P:stable}
For every \(\tau\in\TT(B)\),
\[
f^*(T_\gamma^2\cdot\tau)=f^*(\tau)
\quad\text{and}\quad
i_*(T_\gamma^2\cdot\tau)=i_*(\tau).
\]
In particular, \(T_\gamma^2\) preserves \(\Def_B^A(f)\).
\end{proposition}

\begin{proof}
The pullback map is equivariant with respect to lifts of mapping classes.
By Lemma~\ref{L:twist-lift}, the lift of \(T_\gamma^2\) induces the identity
on \(\TT(A)\).  Therefore
\[
f^*(T_\gamma^2\cdot\tau)=f^*(\tau).
\]
Lemma~\ref{L:forget-trivial} gives
\[
i_*(T_\gamma^2\cdot\tau)=i_*(\tau).
\]
If \(\tau\in\Def_B^A(f)\), then \(f^*(\tau)=i_*(\tau)\), so
\[
f^*(T_\gamma^2\cdot\tau)=i_*(T_\gamma^2\cdot\tau).
\]
Thus \(T_\gamma^2\cdot\tau\in\Def_B^A(f)\).
\end{proof}

This proves the first part of Theorem~\ref{T:main}.

\section{The component-stabilizer computation}

Fix
\[
\tau_0=[\mathrm{id}]\in\Def_B^A(f),
\]
the tautological Teichm\"uller point determined by the given rational map.
Let \(C_0\) be the component of \(\Def_B^A(f)\) containing \(\tau_0\).

Define
\[
S_f:=\{g\in\PMod(\mathbb P^1,B):g\cdot\Def_B^A(f)=\Def_B^A(f)\},
\]
and let
\[
E_f:=\{g\in S_f:g\cdot C_0=C_0\}.
\]
Thus \(E_f\) is the stabilizer in \(S_f\) of the chosen component.

The disconnectedness theorem of Koch--Lodge states that
\(\Def_B^A(f)\) has infinitely many connected components
\cite[Theorem~1.3]{KL-arxiv}.  The proof establishes the following sharper
statement, which is exactly what is needed here.

\begin{theorem}[Koch--Lodge component-stabilizer computation]
\label{T:KL-orbit}
Let \(\gamma\) be the curve separating \(\{c_1,v\}\) from the other three
points of \(B\), and let \(\delta=T_\gamma^2\).  Then
\[
\delta\in S_f,
\]
but
\[
\delta^n\notin E_f\qquad(n\neq0).
\]
More precisely, the cosets
\[
\delta^nE_f\subset S_f,\qquad n\in\mathbb Z,
\]
are pairwise distinct.
\end{theorem}

\begin{proof}
The containment \(\delta\in S_f\) was proved directly in
Proposition~\ref{P:stable}.  The non-containment statement is the explicit
coset computation in the proof of Theorem~1.3 of Koch--Lodge
\cite{KL-arxiv}; see also the published version \cite{KL-published}.
In that proof, \(\delta\) is identified with the point-push of \(v\) around
the distinguished periodic critical point \(c_1\).  The equalizer fundamental
group is then compared with the stabilizer of the deformation space: every
element of the component stabilizer has zero \(c_1\)-winding number, whereas
\(\delta^n\) has \(c_1\)-winding number \(n\).  Hence
\(\delta^nE_f\neq E_f\) for \(n\neq0\), and the cosets
\(\delta^nE_f\) are distinct.

This is stronger than the stated conclusion that the deformation space has
infinitely many components: it is precisely the infinite-coset argument used
to deduce that conclusion.
\end{proof}

\begin{corollary}\label{C:free-orbit}
The components
\[
T_\gamma^{2n}\cdot C_0,\qquad n\in\mathbb Z,
\]
are pairwise distinct.
\end{corollary}

\begin{proof}
Suppose that
\[
T_\gamma^{2n}\cdot C_0=T_\gamma^{2m}\cdot C_0.
\]
Then
\[
T_\gamma^{2(n-m)}\cdot C_0=C_0,
\]
so \(T_\gamma^{2(n-m)}\in E_f\).  Theorem~\ref{T:KL-orbit} forces
\(n-m=0\).  Therefore the components are pairwise distinct.
\end{proof}

This completes the proof of Theorem~\ref{T:main}.

\section{Explicit coordinate description}

Normalize the marked points as
\[
c_1=0,\quad f(c_1)=1,\quad f^{\circ2}(c_1)=\infty,
\quad f^{\circ3}(c_1)=\rho.
\]
Then
\[
B=\{0,1,\infty,\rho,v\}.
\]
The required curve is
\[
\gamma=\partial\Delta,
\]
where \(\Delta\) is a disk containing \(0\) and \(v\) and none of
\(1,\infty,\rho\).  Thus \(\gamma\) separates
\[
\{0,v\}
\quad\text{from}\quad
\{1,\infty,\rho\}.
\]
Its square \(T_\gamma^2\):
\begin{enumerate}
\item lifts under \(f\) to a twist about a curve peripheral relative to \(A\);
\item becomes peripheral after forgetting \(v\);
\item consequently preserves \(\Def_B^A(f)\);
\item permutes its components with a free infinite orbit.
\end{enumerate}

\begin{thebibliography}{99}

\bibitem{KL-arxiv}
S.~Koch and R.~Lodge,
\emph{A disconnected deformation space of rational maps},
arXiv:1602.07378.
HTML version:
\href{https://ar5iv.labs.arxiv.org/html/1602.07378}%
{\texttt{ar5iv.labs.arxiv.org}}.

\bibitem{KL-published}
S.~Koch and R.~Lodge,
\emph{A disconnected deformation space of rational maps},
\emph{Journal of Modern Dynamics} \textbf{11} (2017),
DOI: \href{https://www.aimsciences.org/article/doi/10.3934/jmd.2017016}%
{\texttt{10.3934/jmd.2017016}}.

\bibitem{FKS-arxiv}
T.~Firsova, J.~Kahn, and N.~Selinger,
\emph{On deformation spaces of quadratic rational functions},
arXiv:1606.00723.
HTML version:
\href{https://ar5iv.labs.arxiv.org/html/1606.00723}%
{\texttt{ar5iv.labs.arxiv.org}}.

\bibitem{FKS-pdf}
T.~Firsova, J.~Kahn, and N.~Selinger,
\emph{On deformation spaces of quadratic rational functions},
\href{https://www.math.ksu.edu/~tanyaf/FKS.pdf}{\texttt{math.ksu.edu}}.

\end{thebibliography}

\end{document}